
\section{Methodology}
\label{sec:methodology}

This section describes the proposed prior-aware multimodal framework for long-tailed chest X-ray classification. The model is inspired by CaMCheX, but is adapted to a resource-constrained setting with lightweight image encoding, frozen or precomputed text embeddings, compact multimodal token construction, and prior-aware fusion. Instead of presenting each development stage as a separate model, this section focuses on the final architecture used in the main experiments. Earlier design choices are treated as ablation variants and are discussed in Section~\ref{sec:results}.

\subsection{Problem Formulation}
\label{subsec:method_problem}

This thesis considers multi-label chest X-ray classification with multimodal clinical context. For each current study, the input consists of one or more chest X-ray images, clinical indication text, structured vital signs, and, when available, information from a previous study. The goal is to predict a binary label vector representing the presence or absence of multiple thoracic findings.

Let each sample be represented as:
\begin{equation}
    x_i = (I_i, T_i, V_i, M_i, P_i)
\end{equation}
where \(I_i\) is the set of current chest X-ray views, \(T_i\) is the clinical indication text, \(V_i\) is the vital-sign vector, \(M_i\) is the missingness mask, and \(P_i\) is optional prior-study information. The target label is a multi-hot vector representing the presence of findings:
\begin{equation}
    y_i \in \{0,1\}^{C}
\end{equation}
where \(C\) is the number of pathology labels. The model is trained to learn a parameterized function \(f_\theta\) that produces the class probability predictions:
\begin{equation}
    \hat{y}_i = f_\theta(x_i) \in [0,1]^C
\end{equation}
where each output dimension \(\hat{y}_{ic}\) corresponds to the predicted probability of the \(c\)-th pathology label.

\subsection{Overview of the Proposed Model}
\label{subsec:method_overview}

The proposed model follows a modality-specific encoding and transformer-based fusion \cite{vaswani2017attention} design. First, each input modality is processed by a dedicated encoder. Chest X-ray images are encoded using a lightweight visual backbone. Clinical indication text is encoded using a radiology-domain language model. Vital signs are normalized and projected through a small multilayer perceptron together with their missingness mask. Prior-study information, when available, is encoded separately and used as temporal context.

After encoding, all modality features are projected into a shared latent dimension. Current-study information forms the target stream, while prior-study information forms a memory stream. The fusion module first models interactions among current-study tokens and then allows the current study to attend to prior-study context through cross-attention. This design keeps the current examination as the main prediction target while allowing historical information to support interpretation.

The final fused representation is passed to a multi-label classification head, which outputs one logit for each target finding. The model is trained using an imbalance-aware loss function suitable for long-tailed multi-label classification.

\begin{figure}[htbp]
    \centering
    % \includegraphics[width=0.95\textwidth]{figures/model_architecture.png}
    \fbox{\parbox{0.9\textwidth}{
        \centering
        \vspace{1cm}
        \textbf{Proposed Prior-Aware Multimodal Architecture Diagram}
        \vspace{0.5cm}
        
        \begin{tabular}{ccc}
            \textbf{Target Stream (Current Study)} & & \textbf{Memory Stream (Prior Study)} \\
            \hline
            Current X-Ray Views (up to 4) & & Prior X-Ray Views (up to 4) \\
            $\downarrow$ & & $\downarrow$ \\
            Timm Image Encoder (ConvNeXt) & & Timm Image Encoder (ConvNeXt) \\
            \hline
            Current Clinical Text & & Prior Clinical Report / Text \\
            $\downarrow$ & & $\downarrow$ \\
            BiomedVLP-CXR-BERT & & BiomedVLP-CXR-BERT \\
            \hline
            Structured Vitals \& Missingness Mask & & Prior Vitals \& Missingness Mask \\
            $\downarrow$ & & $\downarrow$ \\
            MLP Projection Layer & & MLP Projection Layer \\
        \end{tabular}
        
        \vspace{0.5cm}
        $\downarrow$ \\
        \textbf{Current Study Fusion (Self-Attention over Target Stream Tokens)} \\
        $\downarrow$ \\
        \textbf{Temporal Cross-Attention (Target Stream attends to Memory Stream Tokens)} \\
        $\downarrow$ \\
        \textbf{Classification Head (MLDecoder with 26 Pathology Queries)} \\
        $\downarrow$ \\
        \textbf{26 Class Probabilities (ASL Loss Optimization)} \\
        \vspace{1cm}
    }}
    \caption{Overview of the proposed prior-aware multimodal chest X-ray classification framework. Current-study image, text, and vital-sign tokens form the target stream, while prior-study information forms the memory stream used through cross-attention.}
    \label{fig:model_architecture}
\end{figure}

\subsection{Data Preprocessing}
\label{subsec:method_preprocessing}

Before feeding the multimodal inputs into their respective encoders, a rigorous preprocessing pipeline is applied to standardise the data formats, handle missing values, and prevent data leakage:

\textbf{Image Preprocessing}: Chest X-ray images are resized to a resolution of $512 \times 512$ pixels. To maintain compatibility with standard pre-trained architectures, single-channel grayscale chest X-rays are duplicated into a 3-channel layout. Alternatively, an optional preprocessing channel mode creates a 3-channel composite combining the raw image, a Contrast Limited Adaptive Histogram Equalization (CLAHE) enhanced image \cite{zuiderveld1994clahe}, and a custom edge-detection channel to emphasize structural details. Pixel values are z-score normalized using the mean and standard deviation computed over the training subset.

\textbf{View Handling}: Studies with multiple radiographic views are limited to a maximum of $4$ views. Views are classified based on their metadata label (\texttt{ViewPosition}): AP, PA, or FRONTAL are mapped to a frontal code ($1$), LATERAL or LL are mapped to a lateral code ($2$), and any other view is coded as $0$. Studies with fewer than $4$ views are zero-padded, and a binary view attention mask is generated so that the transformer fusion layer ignores the padding tokens.

\textbf{Text Tokenization}: Clinical indication texts (and history if present) are parsed and concatenated. Text is tokenized using the tokenizer of the domain-specific \texttt{BiomedVLP-CXR-BERT-}\-\texttt{specialized} model \cite{boecking2022biovil}. The text sequences are padded or truncated to a maximum length of 384 tokens.

\textbf{Vital-Sign Normalization and Outlier Handling}: The structured vital signs vector consists of six numerical parameters (temperature, heart rate, respiratory rate, oxygen saturation, systolic blood pressure, diastolic blood pressure) plus one demographic feature (gender). To handle extreme outliers, values are clipped to physiologically plausible ranges (e.g., temperatures between $85^\circ\text{F}$ and $110^\circ\text{F}$, oxygen saturation capped at $100\%$). Z-score normalization is then performed. For each patient study \(i\) and vital variable \(j\), the normalized value \(\tilde{v}_{ij}\) is formulated as:
\begin{equation}
    \tilde{v}_{ij} =
    \begin{cases}
        \frac{v_{ij} - \mu_j}{\sigma_j}, & \text{if } v_{ij} \text{ is observed}, \\
        0, & \text{if } v_{ij} \text{ is missing},
    \end{cases}
\end{equation}
where \(\mu_j\) and \(\sigma_j\) are the mean and standard deviation of variable \(j\) computed strictly from the training set only to prevent evaluation leakage.

\textbf{Missing Value Imputation}: Missing vital signs are imputed with a neutral z-score value of \(0.0\). To inform the model of data availability, a binary missingness mask \(m_{ij}\) is constructed as:
\begin{equation}
    m_{ij} =
    \begin{cases}
        1, & \text{if } v_{ij} \text{ is observed}, \\
        0, & \text{if } v_{ij} \text{ is missing}.
    \end{cases}
\end{equation}
This mask vector \(M_i = [m_{i1}, \dots, m_{iD}]^\top\) is concatenated to the normalized vital-sign vector \(\tilde{V}_i = [\tilde{v}_{i1}, \dots, \tilde{v}_{iD}]^\top\), enabling the subsequent vital-sign projection network to dynamically adjust its weighting based on data availability.

\textbf{Prior-Study Selection}: To compile longitudinal patient history, studies are sorted chronologically per patient. A shift operation is performed to identify the immediate previous study as the historical prior. This chronological constraint prevents future information from being referenced, thereby ensuring strict temporal leakage prevention.

\textbf{Label Construction}: The binary targets correspond to the $26$ long-tailed pathologies defined by the CXR-LT dataset. Label vectors are represented as $26$-dimensional multi-hot arrays.

\subsection{Modality-Specific Encoders}
\label{subsec:method_encoders}

\subsubsection{Chest X-Ray Image Encoder}
\label{subsubsec:method_image_encoder}

The image encoder extracts visual features from the current chest X-ray study. Because chest X-ray studies may contain different view positions, the model handles frontal and lateral views separately. Each available view is passed through a lightweight ConvNeXt-V2 visual backbone \cite{woo2023convnextv2} to obtain spatial feature maps. These feature maps are then projected to the transformer latent dimension using a convolutional projection layer. Formally, for a set of image views \(I_i\), the image representation \(H_i^{img}\) is extracted as:
\begin{equation}
    H_i^{img} = E_{img}(I_i) \in \mathbb{R}^{N_{img} \times d}
\end{equation}
where \(E_{img}\) denotes the visual backbone and projection layers, \(N_{img}\) is the total number of visual tokens across all views, and \(d\) is the latent dimension of the transformer.

To preserve spatial information, two-dimensional positional embeddings are added to the projected image features. If a study contains fewer views than the maximum expected number of views, missing view slots are represented using padding tokens and ignored by the attention mask. This allows the model to handle variable-view studies while maintaining a consistent tensor structure.

\subsubsection{Clinical Indication Text Encoder}
\label{subsubsec:method_text_encoder}

Clinical indication text provides information about the reason for the examination and the patient's presenting condition. The text is encoded using a radiology-domain language model, BiomedVLP-CXR-BERT-specialized \cite{boecking2022biovil}. The text sequence is tokenized and truncated or padded to a fixed maximum length. The resulting text representation is obtained from the encoder output and projected into the shared fusion dimension:
\begin{equation}
    h_i^{txt} = W_t E_{txt}(T_i) \in \mathbb{R}^{1 \times d}
\end{equation}
where \(E_{txt}\) represents the pre-trained text language model, \(T_i\) is the input indication text, and \(W_t\) is a learnable projection matrix that maps the language model's hidden representation into the shared latent space.

To reduce computational cost, the text encoder is frozen during training. In some experiments, text embeddings are precomputed before training so that the multimodal model does not need to repeatedly run the language model. This design choice makes the experiments feasible under limited GPU resources, although it also means that the text encoder is not adapted end-to-end to the final classification objective.

\subsubsection{Vital Signs Encoder}
\label{subsubsec:method_vitals_encoder}

Structured vital signs provide numerical clinical context related to the patient's physiological state. The vital-sign vector includes variables such as temperature, heart rate, oxygen saturation, respiratory rate, systolic blood pressure, and diastolic blood pressure. Each variable is normalized using statistics computed from the training set.

Because clinical data often contains missing values, the model also uses a missingness mask. The normalized vital-sign vector and missingness mask are concatenated and passed through a small multilayer perceptron to produce a vital-sign token in the shared latent space:
\begin{equation}
    h_i^{vit} = \mathrm{MLP}_v([\tilde{V}_i; M_i]) \in \mathbb{R}^{1 \times d}
\end{equation}
where \([\cdot ; \cdot]\) denotes vector concatenation, and \(\mathrm{MLP}_v\) is a multi-layer perceptron projection network mapping the concatenated physiological vector to the latent dimension \(d\).

\subsection{Multimodal Token Construction}
\label{subsec:method_token_construction}

After modality-specific encoding, the model converts all inputs into tokens for transformer-based fusion. The current-study token sequence \(X_i^{cur}\) is constructed by concatenating the representations from the visual, textual, and physiological encoders:
\begin{equation}
    X_i^{cur} = [H_i^{img}; h_i^{txt}; h_i^{vit}] \in \mathbb{R}^{N_{cur} \times d}
\end{equation}
where \([ \cdot ; \cdot ]\) denotes sequence concatenation along the token dimension, and \(N_{cur} = N_{img} + 2\) is the total number of target tokens.

If prior-study information exists, a corresponding memory stream token sequence \(X_i^{prior}\) is constructed similarly:
\begin{equation}
    X_i^{prior} = [H_i^{prior-img}; h_i^{prior-txt}; h_i^{prior-vit}; e_{\Delta t}] \in \mathbb{R}^{N_{prior} \times d}
\end{equation}
where \(H_i^{prior-img}\), \(h_i^{prior-txt}\), and \(h_i^{prior-vit}\) are the encoded visual, textual, and vital tokens of the historical prior examination, and \(e_{\Delta t}\) is a learned embedding representing the time gap \(\Delta t\) between the current and prior studies.

The use of segment embeddings allows the transformer to distinguish between different sources of information even after all features have been projected into the same latent dimension. Positional embeddings are used for spatial image tokens, while temporal embeddings are used for prior-study tokens. Attention masks are applied to ignore padded views or unavailable modalities.

To make training feasible under limited GPU memory, this thesis utilizes compact token layouts rather than dense spatial maps. For standard self-attention, the computational and memory complexity scales quadratically with the sequence length \cite{vaswani2017attention}:
\begin{equation}
    \mathcal{O}(N^2 d)
\end{equation}
where \(N\) is the number of tokens and \(d\) is the hidden dimension. By projecting spatial visual grids into a reduced set of tokens \(N_{img}\), the attention cost is minimized, allowing multimodal fusion to run within hardware constraints.

\subsection{Prior-Aware Temporal Fusion}
\label{subsec:method_prior_fusion}

Prior-study information is incorporated to reflect the clinical practice of comparing current and previous examinations. The core mechanism of the fusion module is Multi-Head Attention (MHA), following the Transformer attention formulation \cite{vaswani2017attention}. For a query matrix \(Q\), key matrix \(K\), and value matrix \(V\), scaled dot-product attention is defined as:
\begin{equation}
    \mathrm{Attn}(Q,K,V) = \mathrm{softmax}\left(\frac{QK^\top}{\sqrt{d}}\right)V
\end{equation}
where \(d\) is the key projection dimension.

The fusion module uses an asymmetric target-memory design. Current-study tokens form the target stream, while prior-study tokens form the memory stream. The target stream first undergoes self-attention to model interactions among current image, text, and vital-sign tokens:
\begin{equation}
    Z_i^{cur} = \mathrm{MHA}(X_i^{cur}, X_i^{cur}, X_i^{cur})
\end{equation}
It then attends to the prior-study memory stream through cross-attention:
\begin{equation}
    Z_i^{fused} = \mathrm{MHA}(Q=Z_i^{cur}, K=X_i^{prior}, V=X_i^{prior})
\end{equation}
This cross-attention design ensures that the current study remains the query stream that directs the prediction, while the prior study serves as supporting contextual memory.

To represent temporal distances, the time gap \(\Delta t_i\) between studies is mapped into discretized temporal buckets \(b_i\), which are used to retrieve a learned time-gap embedding:
\begin{equation}
    b_i = \mathrm{bucket}(\Delta t_i)
\end{equation}
\begin{equation}
    e_{\Delta t_i} = E_{\Delta t}(b_i)
\end{equation}
where \(E_{\Delta t}\) is a learned embedding lookup table.

When no prior study is available, the memory stream is replaced by a learned no-prior token. During training, stochastic prior dropout is applied with probability \(p_{drop}\). Let \(r_i \in \{0, 1\}\) be a random variable sampled as:
\begin{equation}
    r_i \sim \mathrm{Bernoulli}(1-p_{drop})
\end{equation}
The prior tokens are then scaled during training:
\begin{equation}
    X_i^{prior'} = r_i X_i^{prior}
\end{equation}
This regularization technique reduces the risk of shortcut learning on historical context, encouraging the model to extract features from both current images and prior evidence.

\subsection{Classification Head and Training Objective}
\label{subsec:method_head_loss}

The fused token sequence is passed to a multi-label classification head. Instead of producing a single global pooled representation, the classifier uses learnable label-query embeddings to attend to the fused multimodal features, following the ML-Decoder design \cite{ridnik2023mldecoder}. Let the learnable label queries for the \(C\) target findings be:
\begin{equation}
    Q^{label} = [q_1, q_2, \dots, q_C] \in \mathbb{R}^{C \times d}
\end{equation}
Each label query attends to the fused multimodal token sequence \(Z_i^{fused}\) using multi-head attention:
\begin{equation}
    G_i = \mathrm{MHA}(Q^{label}, Z_i^{fused}, Z_i^{fused}) \in \mathbb{R}^{C \times d}
\end{equation}
The resulting query representations are projected to produce the classification logits:
\begin{equation}
    z_i = W_c G_i + b_c \in \mathbb{R}^{C}
\end{equation}
where \(W_c\) and \(b_c\) are learnable projection weights and biases. The predicted probability for class \(c\) is then computed using the sigmoid function:
\begin{equation}
    \hat{y}_{ic} = \sigma(z_{ic}) = \frac{1}{1 + e^{-z_{ic}}}
\end{equation}

For training, the standard multi-label Binary Cross-Entropy (BCE) loss is defined as:
\begin{equation}
    \mathcal{L}_{BCE} = - \frac{1}{C} \sum_{c=1}^{C} \left[ y_{ic}\log(\hat{y}_{ic}) + (1-y_{ic})\log(1-\hat{y}_{ic}) \right]
\end{equation}
where \(y_{ic}\) is the true binary label.

However, because the dataset is highly imbalanced, this thesis utilizes Asymmetric Loss (ASL) as the actual training objective \cite{ridnik2021asymmetric}. ASL down-weights easy negative examples more aggressively than positive examples. For a single sample, the loss is formulated as:
\begin{equation}
    \mathcal{L}_{ASL} = - \sum_{c=1}^{C} \left[ y_c (1-p_c)^{\gamma_+} \log(p_c) + (1-y_c) p_{mc}^{\gamma_-} \log(1-p_{mc}) \right]
\end{equation}
where \(p_c = \hat{y}_{ic}\) is the predicted probability of class \(c\). To further suppress gradients from easy negatives, the negative probability \(p_{mc}\) is clipped using a threshold \(m \ge 0\):
\begin{equation}
    p_{mc} = \max(p_c - m, 0)
\end{equation}
The shifting hyperparameters \(\gamma_+\) and \(\gamma_-\) control the focusing strength for positive and negative samples, respectively, with \(\gamma_- > \gamma_+\) to focus learning on positive cases while ignoring easy negative background samples.

\subsection{Ablation Study Design}
\label{subsec:method_ablation_design}

Although the main methodology focuses on the final proposed model, several ablation variants are used to understand the contribution of different architectural choices. These ablations are not presented as separate main models, but as controlled comparisons in the experimental analysis.

The ablation studies examine the effect of image-only versus multimodal input, compact versus spatially expanded token layouts, the use of prior-study information, time-gap embeddings, prior dropout, and imbalance-aware loss functions. These comparisons help determine which components contribute most to performance and which design choices are most important under limited computational resources.
